% GENERATED FILE. Edit canonical lessons, then run npm run generate:books.
% canonical-source-hash: fnv1a64:06365119f1318ef4
% canonical-lessons: GE-C297-aufwachsen, GE-C297-einschulung, GE-C297-abitur, GE-C297-abschluss, GE-C297-umzug

\chapter{Aufwachsen}
\label{ch:ge-aufwachsen}
\hlchaptermodality{297}

\begin{chapteropening}
\textbf{By the end of this chapter:} \emph{I can tell how someone grew up: to grow up, starting school, the school-leaving exam, a degree and a move.}

Complete the last lesson of chapter 297: I can tell how someone grew up: to grow up, starting school, the school-leaving exam, a degree and a move.
\end{chapteropening}

\section[aufwachsen]{aufwachsen — to grow up}

\label{lesson:GE-C297-aufwachsen}



\begin{quote}
Before the new one: say the German for to baptise, then the German for the christening.
\end{quote}

\subsection*{You'll want to know: aufwachsen}
\textbf{aufwachsen} — \textquotedblleft{}to grow up\textquotedblright{}.

\begin{quote}
Ich bin auf dem Land aufgewachsen.
\end{quote}

\begin{quote}
\emph{I grew up in the country.}
\end{quote}

Separable, strong, with \emph{sein}: \emph{wuchs auf}, \emph{ist aufgewachsen}.

\begin{cousinweb}
\emph{aufwachsen} joins \emph{auf} \textquotedblleft{}up\textquotedblright{} and \emph{wachsen} \textquotedblleft{}to grow\textquotedblright{}, from Proto-Germanic \emph{wahsijaną} (EtymDB), a cousin of English \emph{wax}.
\end{cousinweb}

\begin{grammarlens}[title={What you've built}]
One more word for telling, explaining and describing at B1.
\end{grammarlens}

\subsection*{Your turn}
\begin{itemize}
\raggedright
  \item \emph{Say it:} \emph{aufwachsen}
  \item \emph{Say it:} \emph{aufwachsen}, once more
  \item \emph{Say it:} a sentence of your own with \emph{aufwachsen}
  \item \emph{Recall:} say the German for to baptise, then the German for the christening, then say \emph{aufwachsen} again
\end{itemize}

\begin{tcolorbox}[breakable,colback=teal!4,colframe=teal!35!black,title={Before you move on}]
What does aufwachsen mean? (\textquotedblleft{}To grow up\textquotedblright{}.) Say it once more.
\end{tcolorbox}
\section[die Einschulung]{die Einschulung — starting school, the first day at school}

\label{lesson:GE-C297-einschulung}



\begin{quote}
Before the new one: say the German for the christening, then the German for to grow up.
\end{quote}

\subsection*{You'll want to know: die Einschulung}
\textbf{die Einschulung} — \textquotedblleft{}starting school, the first day at school\textquotedblright{}.

\begin{quote}
Zur Einschulung bekam sie eine große Schultüte.
\end{quote}

\begin{quote}
\emph{For her first day at school she got a big school cone.}
\end{quote}

A family occasion: the child gets a \emph{Schultüte}, a big paper cone full of sweets.

\begin{cousinweb}
\emph{Einschulung} is built on \emph{Schule} \textquotedblleft{}school\textquotedblright{}, from Old High German \emph{scuola}, which came from Latin \emph{schola} (EtymDB): being taken into school.
\end{cousinweb}

\begin{grammarlens}[title={What you've built}]
One more word for telling, explaining and describing at B1.
\end{grammarlens}

\subsection*{Your turn}
\begin{itemize}
\raggedright
  \item \emph{Say it:} \emph{die Einschulung}
  \item \emph{Say it:} \emph{die Einschulung}, once more
  \item \emph{Say it:} a sentence of your own with \emph{die Einschulung}
  \item \emph{Recall:} say the German for the christening, then the German for to grow up, then say \emph{die Einschulung} again
\end{itemize}

\begin{tcolorbox}[breakable,colback=teal!4,colframe=teal!35!black,title={Before you move on}]
What does die Einschulung mean? (\textquotedblleft{}Starting school, the first day at school\textquotedblright{}.) Say it once more.
\end{tcolorbox}
\section[das Abitur]{das Abitur — the school-leaving exam}

\label{lesson:GE-C297-abitur}



\begin{quote}
Before the new one: say the German for to grow up, then the German for starting school.
\end{quote}

\subsection*{You'll want to know: das Abitur}
\textbf{das Abitur} — \textquotedblleft{}the school-leaving exam\textquotedblright{}.

\begin{quote}
Nach dem Abitur machte er ein Jahr Pause.
\end{quote}

\begin{quote}
\emph{After his school-leaving exam he took a year off.}
\end{quote}

Short: \emph{das Abi}. In Austria and Switzerland: \emph{die Matura}.

\begin{cousinweb}
\emph{Abitur} is a learned formation from Latin \emph{abitūrus} \textquotedblleft{}about to go away\textquotedblright{}, from \emph{abīre} \textquotedblleft{}to go away\textquotedblright{} (Lewis \& Short): the exam of those about to leave school.
\end{cousinweb}

\begin{grammarlens}[title={What you've built}]
One more word for telling, explaining and describing at B1.
\end{grammarlens}

\subsection*{Your turn}
\begin{itemize}
\raggedright
  \item \emph{Say it:} \emph{das Abitur}
  \item \emph{Say it:} \emph{das Abitur}, once more
  \item \emph{Say it:} a sentence of your own with \emph{das Abitur}
  \item \emph{Recall:} say the German for to grow up, then the German for starting school, then say \emph{das Abitur} again
\end{itemize}

\begin{tcolorbox}[breakable,colback=teal!4,colframe=teal!35!black,title={Before you move on}]
What does das Abitur mean? (\textquotedblleft{}The school-leaving exam\textquotedblright{}.) Say it once more.
\end{tcolorbox}
\section[der Abschluss]{der Abschluss — the degree, the qualification; the end}

\label{lesson:GE-C297-abschluss}



\begin{quote}
Before the new one: say the German for starting school, then the German for the school-leaving exam.
\end{quote}

\subsection*{You'll want to know: der Abschluss}
\textbf{der Abschluss} — \textquotedblleft{}the degree, the qualification; the end\textquotedblright{}.

\begin{quote}
Sie machte ihren Abschluss an der Universität.
\end{quote}

\begin{quote}
\emph{She got her degree at the university.}
\end{quote}

Plural: \emph{die Abschlüsse}. Built on \emph{abschließen}.

\begin{cousinweb}
\emph{Abschluss} is built on \emph{Schluss} \textquotedblleft{}end\textquotedblright{}, from \emph{schließen} \textquotedblleft{}to close\textquotedblright{}, Old High German \emph{sliozan} (EtymDB): the closing of a course.
\end{cousinweb}

\begin{grammarlens}[title={What you've built}]
One more word for telling, explaining and describing at B1.
\end{grammarlens}

\subsection*{Your turn}
\begin{itemize}
\raggedright
  \item \emph{Say it:} \emph{der Abschluss}
  \item \emph{Say it:} \emph{der Abschluss}, once more
  \item \emph{Say it:} a sentence of your own with \emph{der Abschluss}
  \item \emph{Recall:} say the German for starting school, then the German for the school-leaving exam, then say \emph{der Abschluss} again
\end{itemize}

\begin{tcolorbox}[breakable,colback=teal!4,colframe=teal!35!black,title={Before you move on}]
What does der Abschluss mean? (\textquotedblleft{}The degree, the qualification; the end\textquotedblright{}.) Say it once more.
\end{tcolorbox}
\section[der Umzug]{der Umzug — the move (to a new home); the parade}

\label{lesson:GE-C297-umzug}



\begin{quote}
Before the new one: say the German for the school-leaving exam, then the German for the degree.
\end{quote}

\subsection*{You'll want to know: der Umzug}
\textbf{der Umzug} — \textquotedblleft{}the move (to a new home); the parade\textquotedblright{}.

\begin{quote}
Der Umzug in die neue Wohnung dauerte zwei Tage.
\end{quote}

\begin{quote}
\emph{The move to the new flat took two days.}
\end{quote}

Plural: \emph{die Umzüge}. Built on \emph{umziehen}. At a festival, a parade.

\begin{cousinweb}
\emph{Umzug} is built on \emph{Zug} \textquotedblleft{}a pull, a move\textquotedblright{}, from \emph{ziehen} \textquotedblleft{}to pull\textquotedblright{}, Proto-Germanic \emph{teuhaną} (EtymDB), with \emph{um} \textquotedblleft{}around\textquotedblright{}.
\end{cousinweb}

\begin{grammarlens}[title={What you've built}]
One more word for telling, explaining and describing at B1.
\end{grammarlens}

\subsection*{Your turn}
\begin{itemize}
\raggedright
  \item \emph{Say it:} \emph{der Umzug}
  \item \emph{Say it:} \emph{der Umzug}, once more
  \item \emph{Say it:} a sentence of your own with \emph{der Umzug}
  \item \emph{Recall:} say the German for the school-leaving exam, then the German for the degree, then say \emph{der Umzug} again
\end{itemize}

\begin{tcolorbox}[breakable,colback=teal!4,colframe=teal!35!black,title={Before you move on}]
What does der Umzug mean? (\textquotedblleft{}The move (to a new home); the parade\textquotedblright{}.) Say it once more.
\end{tcolorbox}
